\documentclass[11pt,a4paper]{article}
\usepackage[T2A]{fontenc}
\usepackage[utf8]{inputenc}
\usepackage[russian,english]{babel}
\usepackage{amsmath,amssymb,amsthm,mathtools}
\usepackage{setspace}
\usepackage{enumitem}
\usepackage[hidelinks]{hyperref}
\onehalfspacing
\newcommand{\head}[1]{\par\medskip\noindent\textbf{\underline{#1}}\quad}
\newcommand{\sheettitle}[3]{% title, subtitle, homework-flag (empty or "Homework")
  \begin{center}
  {\huge\bfseries #1\par}\vspace{1.2em}
  {\Large #2\par}\vspace{0.8em}
  \if\relax\detokenize{#3}\relax\else{\Large #3\par}\vspace{0.8em}\fi
  {\begin{otherlanguage}{russian}\number\day~\ifcase\month\or января\or февраля\or марта\or апреля\or мая\or июня\or июля\or августа\or сентября\or октября\or ноября\or декабря\fi~\number\year~года\end{otherlanguage}\par}
  \end{center}\vspace{0.5em}}
\newcommand{\src}[1]{\unskip\nobreak\hfil\penalty50\hskip1em\hbox{}\nobreak\hfill(#1){\parfillskip=0pt\par}}
\begin{document}
\sheettitle{Euclidean geometry -- 2}{Conics, areas and projective maps}{}

\head{Theoretical minimum.} A conic is seen metrically (focus and directrix),
algebraically (a quadratic equation, classified by linear algebra) and
projectively (a curve fixed by five points, carrying cross-ratios); for areas
one first applies an affine map turning it into a circle, $xy=1$ or $y=x^2$.

\head{Definitions.} For a point $F$, a line $l\not\ni F$ and $e>0$ the
\emph{conic} $\{X: |XF|=e\,d(X,l)\}$ with focus $F$, directrix $l$ and
eccentricity $e$ is an ellipse, a parabola or a hyperbola for $e<1$, $e=1$,
$e>1$; in polar coordinates at $F$, with $\varphi$ measured from the
perpendicular from $F$ to $l$, it is $r=p/(1+e\cos\varphi)$, $p=e\,d(F,l)$.
The \emph{signed area} is $[ABC]=\frac12\det(B-A,C-A)$, positive for
counterclockwise $ABC$. The \emph{cross-ratio} of points $a,b,c,d$ of a line
is $(a,b;c,d)=\frac{(c-a)(d-b)}{(c-b)(d-a)}$; a \emph{projective map} of
$\mathbb R\cup\{\infty\}$ is $x\mapsto\frac{\alpha x+\beta}{\gamma x+\delta}$,
$\alpha\delta\ne\beta\gamma$; a \emph{perspectivity} from a line $l$ to a line
$m$ with centre $O\notin l\cup m$ sends $X$ to $OX\cap m$.

\head{Theorem 1 (focal and reflective properties).} For $e\ne1$ there is a
second focus and directrix, symmetric in the centre, and
$|XF_1|+|XF_2|=2a$ (ellipse), $\bigl||XF_1|-|XF_2|\bigr|=2a$ (hyperbola), with
$c=|OF_i|=ea$ and directrices at distance $a/e$ from the centre $O$. The
tangent at $X$ bisects the external angle $F_1XF_2$ of an ellipse and the
internal one of a hyperbola; for a parabola it bisects the angle between $XF$
and the perpendicular from $X$ to $l$. \emph{Proof idea:} $X$ minimises
$|YF_1|+|YF_2|$ over the tangent, since the other points of the tangent lie
outside; then use Heron's reflection. So light from one focus is reflected
into the other one (parallel to the axis for a parabola).

\head{Theorem 2 (equations, five points, pencils).} A conic is the zero set
of $Q=ax^2+2bxy+cy^2+2dx+2ey+f$; up to affine maps a non-degenerate non-empty
real conic is $x^2+y^2=1$, $x^2-y^2=1$ or $y=x^2$ as $ac-b^2>0$, $<0$, $=0$
(Linear algebra -- 5), and projectively all three coincide. A line not
contained in a conic meets it at most twice; as $Q$ has six coefficients,
through five points, no four collinear, passes exactly one conic. The conics
through four points, no three collinear, form a \emph{pencil}
$\lambda Q_1+\mu Q_2$ containing exactly three line pairs: the two pairs of
opposite sides and the pair of diagonals.

\head{Theorem 3 (areas).} In coordinates
\[
[ABC]=\tfrac12\bigl((x_B-x_A)(y_C-y_A)-(x_C-x_A)(y_B-y_A)\bigr),
\]
so $P\mapsto[PBC]$ is affine and $[PBC]+[PCA]+[PAB]=[ABC]$ for all $P$
(barycentric coordinates); at an interior extremum of $[PAB]$ over a smooth
arc the tangent is parallel to $AB$. An affine map $x\mapsto Mx+v$ multiplies
areas by $\det M$ and preserves area ratios, parallelism, tangency and
centres; all parabolic segments are affinely equivalent, and $y=x^2$ cuts off
over $[a,b]$ a segment of area $(b-a)^3/6$.

\head{Theorem 4 (projective maps of a line).} A perspectivity with centre
$O$ preserves the cross-ratio, since
$(A,B;C,D)=\frac{\sin\angle AOC\,\sin\angle BOD}{\sin\angle BOC\,\sin\angle AOD}$
(oriented angles). Projective maps are exactly the cross-ratio preserving
bijections of $\mathbb R\cup\{\infty\}$; each is determined by the images of
three points, so one with three fixed points is the identity. Compositions of
perspectivities are projective, and every projective map of a line to itself
is a composition of at most three perspectivities (over $\mathbb C$: the
M\"obius maps of Complex analysis -- 7).

\head{Fact 1 (normals of a parabola, Apollonius).} The normal to $x^2=2py$ at
$(t,t^2/2p)$ passes through $(u,v)$ iff $t^3+2p(p-v)t-2p^2u=0$: at most three
normals pass through a point, the abscissas of their feet add up to $0$, and
there are three exactly when $8(v-p)^3>27pu^2$.

\head{Fact 2 (billiards, Birkhoff).} In a smooth strictly convex curve $\Gamma$
an inscribed $n$-gon of maximal perimeter exists by compactness; at each
vertex $X$ the sum $|YX_-|+|YX_+|$ over $Y\in\Gamma$ is maximal at $Y=X$, so
the normal bisects $\angle X_-XX_+$ and the polygon is a closed billiard
trajectory --- the argument of Theorem 1 again.

\medskip\noindent\rule{\linewidth}{0.4pt}

\begin{enumerate}[leftmargin=2.2em,itemsep=0.6em]
\item Derive the polar equation from $|XF|=e\,d(X,l)$ and, for $e\ne1$, the
canonical equation $\frac{x^2}{a^2}\pm\frac{y^2}{b^2}=1$ with $c=ea$ and the
directrices $x=\pm a/e$; prove the focal-sum property of the ellipse and the
reflective properties of Theorem 1, and show that a billiard ball sent
through one focus of an elliptic table passes through the other focus after
every reflection. Then prove Fact 1.

\item Prove the determinant formula of Theorem 3 and the identity
$[PBC]+[PCA]+[PAB]=[ABC]$. Prove that a perspectivity preserves the
cross-ratio, that a projective map of a line is determined by the images of
three points, and deduce that a projective map other than the identity has at
most two fixed points. Finally, prove that through five points of which no
four are collinear there passes exactly one conic.

\item Two different ellipses are given. One focus of the first ellipse
coincides with one focus of the second ellipse. Prove that the ellipses have
at most two points in common. \src{IMC 2008}

\item The normals to the parabola $2y=ax^2$ $(a>0)$ at two different points
$B$ and $B'$ of the parabola meet at a point $A$ of the parabola, different
from $B$ and $B'$; let $m$ be the ordinate of $A$. Prove that $m>\frac4a$.
\src{IMC 2025, shortlist}

\item An ellipse, a parabola and a hyperbola are circumscribed about a
trapezoid. Prove that the line joining the centres of the ellipse and of the
hyperbola is parallel to the axis of the parabola. \src{МФТИ 2016}
\end{enumerate}
\end{document}
